\subsection{Secure \acs{API} Exposure via \acs{CAPIF}}
\label{chap:proposal:exposure:capif}

The \acs{NEF} and CAMARA components described in the preceding sections provide the mechanisms necessary to expose network capabilities to third parties. Together, they provide developers with the interfaces required to interact with the underlying \acs{5G} infrastructure, leveraging it to build network applications. Nevertheless, a standardized mechanism for the secure discovery, exposure, and consumption of these northbound \acsp{API} is yet to be defined.

As introduced at the beginning of this section, these northbound \acsp{API} should be exposed through a common governance mechanism, i.e. \acs{CAPIF}~\cite{ETSI_CAPIF_TS_123_222_V18_4_0}. It introduces a set of roles, namely the \acs{API} Provider, \acs{API} Invoker and \acf{CCF}, along with the procedures for \acs{API} publishing, discovery, and access control. Given its generic design, several works in the research community have treated \acs{CAPIF} as a general-purpose \acs{API} gateway with authentication and authorization capabilities, leveraging it to protect both \acs{NEF} and CAMARA northbound \acsp{API}~\citep{PathwayNaaSCAMARA,pushingBoundariesScalable5GCN}.

This dissertation follows the same interpretation, using \acs{CAPIF} as a unified secure exposure layer that consolidates common supporting capabilities into a single component, for all northbound \acsp{API} in the testbed. Under this model, (i) \acs{CCF} acts as the central control entity of the framework; (ii) both the \acs{NEF} and CAMARA \acsp{API} act as \acs{API} providers exposing interfaces to be consumed by network applications, which operate as external \acs{API} invokers; and (iii) each CAMARA \acs{API} additionally assumes the role of an internal \acs{API} invoker, interacting with \acs{NEF}'s \acs{API} to deliver its services from inside the network's trust domain.

{\spaceskip=\fontdimen2\font plus 0.5em minus \fontdimen4\font
The procedures defined in the \acs{CAPIF} architecture are realized as a collection of RESTful \acsp{API} in \acs{3GPP} TS 29.222~\cite{ETSI_CAPIF_TS_129_222_V18_6_0}, namely the \texttt{Discover\allowbreak\_Service}, \texttt{Publish\allowbreak\_Service}, \texttt{Events}, \texttt{API\allowbreak\_Invoker\allowbreak\_Management}, \texttt{Security}, \texttt{Monitoring}, \texttt{Logging\allowbreak\_API\allowbreak\_Invocation}, \texttt{Auditing}, \texttt{Access\allowbreak\_Control\allowbreak\_Policy}, \texttt{API\allowbreak\_Provider\allowbreak\_Management}, and \texttt{Routing\allowbreak\_Info} \acsp{API}.}
In the testbed, the exposure of each northbound \acs{API} would follow these procedures. Before any \acs{API} becomes available, each provider should register with the \acs{CCF} and publish its service \acs{API} information through the \acs{CAPIFAPF}. Each network application would then onboard with the \acs{CCF} to obtain its security credentials, and use the discovery interface of the \acs{CCF} to identify the exposed \acsp{API}, their endpoints, and the authorization policies governing their use. Each invocation would subsequently reach the \acs{CAPIFAEF}, which authenticates the invoker, validates the authorization issued by the \acs{CCF}, and logs the invocation for auditing. Through these onboarding and authorization policies, the set of \acsp{API} available to each network application can also be restricted~\cite{ETSI_CAPIF_TS_123_222_V18_4_0}.

As a result, rather than requiring separate onboarding, discovery, and authentication procedures for each \acs{API} suite, a single \acs{CAPIF} integration point consolidates these concerns regardless of whether the underlying \acs{API} is a \acs{3GPP} \acs{NEF} interface or a CAMARA abstraction. This significantly reduces integration complexity and aligns with the testbed's objective of providing a developer-friendly environment for the rapid development and validation of network applications.
